\section{HMT reading duality: duality \texorpdfstring{\(T\)}{T}, block \texorpdfstring{\(243\)}{243}, ternary Golay code, and mass routes}
\label{sec:sucesor83-f083-dualidad-t}

The \(T\) duality of strings is realized as the continuous projection of a
prior HMT duality: the exchange between visible value and boundary
memory. Every integer admits the reading
\[
n=9Q+r,
\]
where \(r\) is the visible value of the nonadic wheel and \(Q\) preserves
boundary-crossing memory. The same structure appears in lifted APP,
in the decomposition \(1215=54+9\cdot129\), in the stratified dimensional
unit \(1000=729+243+27+1\), and in the reading of mass as a route with
visible advance and counterangular winding. This duality links the block
\(243\), the perfect ternary ball \(V_3(11,2)=243\), the punctured ternary
Golay code \([11,6,5]_3\), the channels of
\(\pi,e,\varphi,\alpha\), the hinge \(744\), the dimension
\(11=6+5=3+8\), and the involution
\[
(r,Q;R)\longmapsto(Q,r;R^{-1}).
\]
In string-theoretic language, the visible value is published as a vibrational mode or
momentum, whereas boundary memory is published as winding.
In Dimensional Mechanics these two modes appear as
\(n_AA_{\mathrm{rad}}\) and \((s_C/6)C^*_{\mathrm{rad}}\): visible advance and
counterangular memory.

HMT does not take the continuum as a primary object. It starts from a discrete cell
with boundary memory. The real line, Moonshine, and dimensional physics are
different readings of that same cell. Consequently, the chain
\[
\text{lifted APP}
\longrightarrow\text{value and memory}
\longrightarrow\text{duality }T
\longrightarrow243/\text{Golay}
\longrightarrow\text{mass routes}
\]
does not import duality from string theory: it shows that its structure
exists in HMT before continuous geometrization.

\subsection{Visible residue and boundary quotient}

For every positive integer \(n\), let be defined
\[
r_9(n)=1+((n-1)\bmod 9),
\qquad
Q_9(n)=\frac{n-r_9(n)}{9}.
\]
Then
\[
n=9Q_9(n)+r_9(n).
\]
The ordinary reading preserves the residue. The enriched reading preserves the complete pair.
\[
n\longmapsto(r,Q).
\]
This pair is the minimal structure of a compact reading: one component is
published inside the wheel and the other preserves how many times its
boundary was traversed.

\begin{principle}[Duality between visible value and boundary memory]
The HMT cell admits two conjugate readings. The first publishes \(r\) as
phase, value, or vibration; the second publishes \(Q\) as memory,
winding, or number of boundary crossings.
\end{principle}

\subsection{T-duality involution}

Let \(R>0\). The functional
\[
E_R(r,Q)=\frac{r^2}{R^2}+Q^2R^2
\]
satisfies
\[
E_R(r,Q)=E_{1/R}(Q,r),
\]
for
\[
E_{1/R}(Q,r)
=\frac{Q^2}{(1/R)^2}+r^2(1/R)^2
=Q^2R^2+\frac{r^2}{R^2}.
\]
Therefore,
\[
(r,Q;R)\longleftrightarrow(Q,r;R^{-1}).
\]
The continuous string expression contains terms of the form
\[
\frac{n^2}{R^2}+\frac{w^2R^2}{\alpha'^2}.
\]
Before that continuous realization, HMT supplies the visible residue, boundary
memory, and scale inversion. The duality \(T\) is the geometric
projection of this reading involution.

\subsection{APP raised: visible modes and memory modes}

For the sum and product over the same support one writes
\[
i+j=9Q^+_{ij}+R^+_{ij},
\qquad
ij=9Q^\times_{ij}+R^\times_{ij}.
\]
The raised APP is
\[
\mathcal A_9^\sharp=(R^+,Q^+;R^\times,Q^\times),
\]
where \(R^+,R^\times\) are the visible modes and
\(Q^+,Q^\times\) the memory modes. The complete discrepancy is
\[
ij-(i+j)
=(R^\times_{ij}-R^+_{ij})
+9(Q^\times_{ij}-Q^+_{ij}).
\]
When summing over \(1\leq i,j\leq9\),
\[
\sum R^+=405,
\quad
\sum R^\times=459,
\quad
\sum Q^+=45,
\quad
\sum Q^\times=174.
\]
From here, one obtains
\[
\Delta_{\mathrm{visible}}=459-405=54,
\qquad
\Delta_{\mathrm{memoria}}=174-45=129,
\]
and
\[
\sum ij-\sum(i+j)=2025-810=1215=54+9\cdot129.
\]
Normalization by the wheel produces
\[
\frac{1215}{9}=135=6+129=120+15.
\]
Thus, \(54\) is the visible defect, \(129\) the memory defect,
\(135\) the normalized total defect, \(120\) the constitutive octadic channel,
and \(15\) the internal pair of the hexad.

\subsection{Stratified completion of the unit}

The Extreme and Mean Holographic Ratio possesses the cubic completion
\[
10^3=9^3+3\cdot9^2+3\cdot9+1,
\]
that is,
\[
1000=729+243+27+1.
\]
The four summands represent, respectively, the discrete volume of
dimension \(3\), the boundary of codimension \(1\), the boundary of
codimension \(2\), and the return of codimension \(3\). The active corona
is refined as
\[
270=243+27=3\cdot9^2+3\cdot9.
\]
The unit is thus decomposed into interior, stratified active boundary, and
return. The passage \(9\to10\) opens a complete dimensional layer, and the boundary
of one level becomes a propagation channel of the next level.

\subsection{The block \texorpdfstring{\(243\)}{243} as boundary and perfect ternary ball}

The same block is obtained as a ternary Hamming ball of radius \(2\) and
length \(11\):
\[
V_3(11,2)
=\sum_{k=0}^{2}\binom{11}{k}2^k
=1+22+220
=243.
\]
The punctured perfect ternary Golay code has parameters
\[
G_{11}=[11,6,5]_3,
\qquad
|G_{11}|=3^6=729,
\]
and satisfies
\[
729\cdot243=3^{11}.
\]
Therefore, \(G_{11}\) tiles \(\mathbb F_3^{11}\) by ternary balls of
radius \(2\). The number \(243\) is simultaneously
\[
3\cdot9^2,
\qquad
V_3(11,2),
\qquad
3^5,
\qquad
1+22+220.
\]
This identity joins dimensional boundary, perfect correction, dimension
\(11\), Witt--Golay incidence and the structural reading of M-theory.

\subsection{Internal decomposition of the block \texorpdfstring{\(243\)}{243}}

The ball admits the decomposition
\[
243=1+22+90+120+10.
\]
Here \(1\) is the center and return; \(22=2\cdot11\), the first Hamming shell
and the electronic reading of the block; \(90=6\binom62\), the hexadic channel;
\(120=8\binom62\), the octadic channel; and \(10=9+1\), closure of the
monodromy \(9\to1\). Thus, the return, component \(22\), channels \(90/120\), and
monodromy are contained in the same perfect ball.

\subsection{Constant channels derived from the block \texorpdfstring{\(243\)}{243}}

With \(B=243\), one obtains the channels
\[
C_\alpha=3B=729,
\qquad
C_e=B+27+1=271,
\]
\[
C_\pi=B+72=315,
\qquad
C_\varphi=B-81-1=161.
\]
The channel of \(\alpha\) is the triple block \(243\); that of \(e\) is the block
\(243\) with the boundary \(27\) and the return; that of \(\pi\) is the block with the
channel \(72\); and that of \(\varphi\) is the block after removing the
subsquare \(81\) and the return.
The hinge is reconstructed as
\[
C_\pi+C_e+C_\varphi-3
=315+271+161-3
=744
=729+15.
\]
The reduced trace of the three characters therefore combines the triple block
\(243\) with the internal pair of the hexad.

\subsection{Dimension \texorpdfstring{\(11\)}{11} and compatible readings}

The punctured perfect ternary code produces
\[
11=6+5,
\]
where \(6\) is the dimension of \(G_{11}\) and \(5\) is the correction depth \(3^5=243\). The same dimension decomposes as
\[
11=3+8,
\]
where \(3\) corresponds to TRIT and \(8\) to the internal octad. Taken together,
\[
11=6+5=3+8,
\qquad
12=11+1.
\]
The dimension reading \(11\) thus coordinates code and correction, TRIT and octad, and the identity \(729\cdot243=3^{11}\).

\subsection{Dimensional realization of duality}

The local action of a route is expressed as
\[
E_0(P)=n_A(P)A_{\mathrm{rad}}
+\frac{s_C(P)}{6}C^*_{\mathrm{rad}}.
\]
The first term is the visible advance or vibration mode; the second is the
boundary or contra-angular winding mode. The local mass is
\[
m_{\mathrm{HMT}}(P)=m_e\exp(E_0(P)).
\]
The refined tower adds the boundary memory:
\[
E(P)=E_0(P)+k(P)\Delta_4
+\nu_{120}(P)s_{120}
+\nu_{270}(P)\zeta_{270},
\]
where \(\Delta_4\) is the minimal torsion, \(s_{120}\) the
hexad--octad selector, and
\[
\zeta_{270}=135\Delta_4^2
\]
the quadratic memory of the crown. The abstract duality of route is
\[
E_{A,C}(n,w)=nA_{\mathrm{rad}}+wC^*_{\mathrm{rad}},
\qquad
E_{A,C}(n,w)=E_{C,A}(w,n).
\]
In the physical sector, \(A\) and \(C^*\) are polarized, and the genealogical register selects
the realization. Mass is the exponential of a closure route with visible advance,
boundary winding, and torsion.

\subsection{Compatibility between leaves and closure of the genealogical record}

In a closed string, compatibility between the left- and right-moving modes
is expressed schematically by
\[
N_L-N_R=nw.
\]
In the HMT reading, the left-hand member measures the difference of excitation
between sheets, whereas the right-hand member measures the coupling between the visible phase and
boundary memory. The adjustment of levels thus constitutes a closure of the
bidirectional genealogical record. In the readings \(CP\), \(T\), and \(CPT\), an asymmetry
of orientation is compensated through bidirectional inversion in order to close the
record.

\subsection{T-Duality, Modularity and J Function}

In the complex torus,
\[
q=e^{2\pi i\tau},
\]
and the modular transformation
\[
\tau\longmapsto-\frac1\tau
\]
interchanges its cycles. The duality \(T\) interchanges large and small radii; both realizations publish the interchange between the visible cycle and the boundary cycle. The modular function \(J\), the Sommerfeld action, and the duality \(T\) constitute distinct closures of the same primary boundary structure with memory.

\subsection{Dimensional Consequences for Photon, Electron, Mass, and Statistics}

The observable photon is associated with the central TPK mode:
\[
L_{\mathrm{TPK}}=3(I-P_0),
\qquad
\operatorname{Spec}(L_{\mathrm{TPK}})=\{0,3,3\}.
\]
The zero mode is an extended null propagation with no compactified memory
of its own. The electron appears as the compact closure of the conjugate
zero modes,
\[
e^-=\text{compact closure of }P_++P_-,
\]
and constitutes the minimal physical winding with stable memory. Mass
measures the cost of stabilizing that compact memory. Pauli
exclusion is interpreted as saturation of the genealogical register: two identical fermions do not
occupy the same boundary because their winding closure is already saturated.
Bosons publish shareable modes of unfolded memory, whereas
fermions publish compactified memory.

\subsection{Verifiable identities of the construction}

The construction satisfies jointly the identities
\begin{enumerate}
\item \(1000=729+243+27+1\);
\item \(V_3(11,2)=1+22+220=243\);
\item \(729\cdot243=3^{11}\);
\item \(243=1+22+90+120+10\);
\item \(C_\alpha,C_e,C_\pi,C_\varphi=729,271,315,161\);
\item \(315+271+161-3=744\);
\item \(1215=54+9\cdot129\);
\item \(E_R(r,Q)=E_{1/R}(Q,r)\);
\item \(11=6+5=3+8\), \(12=11+1\), \(26=2+24\) y \(10=2+8\).
\end{enumerate}

The duality \(T\) is therefore recognized as the geometric realization of the
exchange between residue and quotient. Winding is not introduced
first as an external property of a compact string, but as boundary memory of
the decomposition \(n=9Q+r\). The block \(243\) brings this
duality together with ternary Golay, dimension \(11\), channels \(90/120\), generation
of constants, mass routes, and M-theory. In its most general form, HMT reads
mass, string, and dimension as publications of a single boundary with
memory.
